%! Rates of Change in Polar Functions — Unit 3, Lesson 3.15 — Guided Notes (Answer Key)
\documentclass[10pt]{article}
\usepackage{precalculus-article}
\usepackage{precalculus-key}   % pulls in -boxes; do NOT also load -boxes
\newcommand{\TallMath}[1]{$\displaystyle #1\rule[-1.4em]{0pt}{3.2em}$}
\newcommand{\CourseName}{AP Precalculus}
\newcommand{\SchoolYear}{2026--2027}
\newcommand{\MeetingLength}{55 minutes}
\newcommand{\UnitNumberName}{Unit 3: Trigonometric and Polar Functions \quad}
\newcommand{\LessonNumberName}{Lesson 3.15: Rates of Change in Polar Functions}

\begin{document}

\pageheader{Unit 3, Lesson 3.15}{Guided Notes --- Answer Key}
\namedateperiod

\vspace{0.10in}

%----------------------------------------------------------------------
% Objectives
%----------------------------------------------------------------------
\begin{objectivebox}
\textbf{I will be able to\ldots}
\begin{enumerate}[leftmargin=1.5em, itemsep=4pt]
  \item \textbf{LO 3.15.A} --- Compute the average rate of change of a polar
    function $r = f(\theta)$ over an interval and interpret it as how quickly
    the radius changes per radian.
    \ans{Use $(f(\theta_2)-f(\theta_1))/(\theta_2-\theta_1)$; positive $\Rightarrow$ $r$ growing, negative $\Rightarrow$ $r$ shrinking.}
  \item \textbf{LO 3.15.A} --- Classify intervals as P\&I, P\&D, N\&I, or N\&D
    and describe the corresponding curve behavior.
    \ans{P\&I: away from origin; P\&D: toward origin; N\&I: toward origin (opposite ray); N\&D: away from origin (opposite ray).}
\end{enumerate}
\end{objectivebox}

\vspace{0.08in}

%----------------------------------------------------------------------
% Vocabulary
%----------------------------------------------------------------------
\begin{vocabbox}
\begin{description}[leftmargin=0pt, itemsep=6pt]
  \item[\termblanklong{rate of change of a polar function}]
    For $r = f(\theta)$, the average rate of change over $[\theta_1, \theta_2]$ is
    $\dfrac{f(\theta_2) - f(\theta_1)}{\theta_2 - \theta_1}$. Measures how quickly
    $r$ changes per \ans{radian} of angle change.

  \item[\termblanklong{positive and increasing}]
    $r > 0$ and $r$ is \ans{growing}. The curve moves
    \ans{away} from the origin as $\theta$ increases.

  \item[\termblanklong{positive and decreasing}]
    $r > 0$ and $r$ is \ans{shrinking}. The curve moves
    \ans{toward} the origin as $\theta$ increases.

  \item[\termblanklong{negative and increasing}]
    $r < 0$ and $r$ is becoming \ans{less} negative (closer to zero).
    The curve is traced on the \ans{opposite} ray, moving toward the origin.

  \item[\termblanklong{negative and decreasing}]
    $r < 0$ and $r$ is becoming \ans{more} negative (farther from zero).
    The curve is traced on the opposite ray, moving \ans{away from} the origin.
\end{description}
\end{vocabbox}

\vspace{0.08in}

%----------------------------------------------------------------------
% Hook
%----------------------------------------------------------------------
\begin{hookbox}
\textbf{Entry Question:} The table below shows values of $r = 1 + \sin\theta$.

\vspace{4pt}
\renewcommand{\arraystretch}{1.4}
\begin{tabular}{|c|c|c|c|c|c|c|c|c|}
\hline
$\theta$ & $0$ & $\pi/4$ & $\pi/2$ & $3\pi/4$ & $\pi$ & $5\pi/4$ & $3\pi/2$ & $7\pi/4$ \\
\hline
$r$ & 1 & $\approx 1.71$ & 2 & $\approx 1.71$ & 1 & $\approx 0.29$ & 0 & $\approx 0.29$ \\
\hline
\end{tabular}

\vspace{8pt}
\textbf{Q1.} \ans{$[0, \pi/2]$: $r$ increases from 1 to 2.}

\vspace{6pt}
\textbf{Q2.} \ans{$[\pi/2, 3\pi/2]$: $r$ decreases from 2 to 0.}

\vspace{6pt}
\textbf{Q3.} \ans{$r = 0$ means the point is at the origin (the pole). The curve passes through the origin.}

\end{hookbox}

\vspace{0.08in}

%----------------------------------------------------------------------
% LO 3.15.A — Average Rate of Change
%----------------------------------------------------------------------
\begin{notesbox}{LO 3.15.A --- Average Rate of Change of $r$ --- Key}

\textbf{Formula:}
$$\text{Average Rate of Change} = \frac{f(\theta_2) - f(\theta_1)}{\theta_2 - \theta_1}$$

This tells us: for every additional \ans{radian} that $\theta$ increases,
$r$ changes by approximately \ans{this amount (the ARC value)}.

\vspace{8pt}
\textbf{Worked Example:} $r = 2\sin\theta$

\vspace{4pt}
\renewcommand{\arraystretch}{1.5}
\begin{tabular}{|c|c|c|}
\hline
$\theta$ & $0$ & $\pi/2$ \\
\hline
$r = 2\sin\theta$ & $0$ & $2$ \\
\hline
\end{tabular}

\vspace{8pt}
$$\frac{2 - 0}{\pi/2 - 0} = \frac{4}{\pi}$$

Result $= $ \ans{$4/\pi \approx 1.27$ per radian}

\vspace{6pt}
Interpret: $r$ increases by approximately \ans{$4/\pi \approx 1.27$} per radian on average.

\vspace{8pt}
\textbf{Your turn:} ARC from $\theta = \pi/2$ to $\theta = \pi$ for $r = 2\sin\theta$:

\vspace{6pt}
ARC $= \dfrac{{\color{keyred}\mathbf{0}} - {\color{keyred}\mathbf{2}}}{{\color{keyred}\mathbf{\pi - \pi/2}}} = $ \ans{$-4/\pi \approx -1.27$ per radian}

\vspace{4pt}
$r$ is \ans{decreasing} on $[\pi/2, \pi]$ (negative ARC).

\end{notesbox}

\vspace{0.08in}

%----------------------------------------------------------------------
% LO 3.15.A — Sign and Direction
%----------------------------------------------------------------------
\begin{notesbox}{LO 3.15.A --- Sign and Direction of $r$ --- Key}

\vspace{6pt}
\renewcommand{\arraystretch}{1.7}
\begin{tabular}{|l|l|p{4.5cm}|}
\hline
\textbf{Sign of $r$} & \textbf{Behavior of $r$} & \textbf{What happens in the plane} \\
\hline
$r > 0$ & increasing & \ans{Moves away from origin} \\
\hline
$r > 0$ & decreasing & \ans{Moves toward origin} \\
\hline
$r < 0$ & increasing (less negative) & \ans{Moves toward origin (opposite ray)} \\
\hline
$r < 0$ & decreasing (more negative) & \ans{Moves away from origin (opposite ray)} \\
\hline
\end{tabular}

\vspace{8pt}
\textbf{Key reminder:} When $r < 0$, the point is plotted on the
\ans{opposite} direction from the terminal ray, at distance
$|r|$ from the origin.

\end{notesbox}

\vspace{0.08in}

%----------------------------------------------------------------------
% LO 3.15.A — Intervals from a Table
%----------------------------------------------------------------------
\begin{notesbox}{LO 3.15.A --- Reading Intervals from a Table --- Key}

$r = 1 + \cos\theta$:

\vspace{4pt}
\renewcommand{\arraystretch}{1.4}
\begin{tabular}{|c|c|c|c|c|c|c|c|c|}
\hline
$\theta$ & $0$ & $\pi/4$ & $\pi/2$ & $3\pi/4$ & $\pi$ & $5\pi/4$ & $3\pi/2$ & $7\pi/4$ \\
\hline
$r$ & $2$ & $\approx 1.71$ & $1$ & $\approx 0.29$ & $0$ & $\approx 0.29$ & $1$ & $\approx 1.71$ \\
\hline
\end{tabular}

\vspace{8pt}
\begin{tabular}{|c|p{5cm}|}
\hline
\textbf{Interval} & \textbf{Classification} \\
\hline
$[0, \pi/4]$ & \ans{Positive and decreasing (P\&D)} \\
\hline
$[\pi/4, \pi/2]$ & \ans{Positive and decreasing (P\&D)} \\
\hline
$[\pi/2, 3\pi/4]$ & \ans{Positive and decreasing (P\&D)} \\
\hline
$[3\pi/4, \pi]$ & \ans{Positive and decreasing (P\&D)} \\
\hline
$[\pi, 5\pi/4]$ & \ans{Positive and increasing (P\&I)} \\
\hline
$[5\pi/4, 3\pi/2]$ & \ans{Positive and increasing (P\&I)} \\
\hline
$[3\pi/2, 7\pi/4]$ & \ans{Positive and increasing (P\&I)} \\
\hline
\end{tabular}

\vspace{6pt}
$r = 0$ at $\theta = $ \ans{$\pi$}

\end{notesbox}

\vspace{0.08in}

%----------------------------------------------------------------------
% Practice
%----------------------------------------------------------------------
\begin{practicebox}

\textbf{Guided Practice 1.}\quad $r = 3\cos\theta$:

\vspace{4pt}
\renewcommand{\arraystretch}{1.5}
\begin{tabular}{|c|c|c|c|c|c|}
\hline
$\theta$ & $0$ & $\pi/4$ & $\pi/2$ & $3\pi/4$ & $\pi$ \\
\hline
$r$ & $3$ & $\approx 2.12$ & $0$ & $\approx -2.12$ & $-3$ \\
\hline
\end{tabular}

\vspace{6pt}
(a)\quad $[0, \pi/4]$: \ans{P\&D} \quad
$[\pi/4, \pi/2]$: \ans{P\&D} \quad
$[\pi/2, 3\pi/4]$: \ans{N\&D} \quad
$[3\pi/4, \pi]$: \ans{N\&D}

\vspace{0.10in}
\textbf{Guided Practice 2.}

\vspace{4pt}
(a)\quad ARC $= \dfrac{{\color{keyred}\mathbf{0}} - {\color{keyred}\mathbf{3}}}{\pi/2} = $ \ans{$-6/\pi \approx -1.91$ per radian}

\vspace{6pt}
(b)\quad ARC $= \dfrac{-3 - 0}{\pi/2} = $ \ans{$-6/\pi \approx -1.91$ per radian}

\vspace{6pt}
(c)\quad \ans{On $[\pi/2, \pi]$, $r < 0$ and decreasing (N\&D), so the curve moves away from the origin on the opposite ray.}

\end{practicebox}

\end{document}
